% OWNER: theory-b
% STATUS: review
% SOURCES: notes/SPIS_TRESCI_SZCZEGOLOWY.md §3.5--3.6 / most osobowości;
%          notes/HIPOTEZY_REWIZJA.md (RQ4); GLOSSARY.md; appendix-c-mappings;
%          research-mappings.ts / facet-derivation.ts (hipotezy projektowe)
\chapter{Osobowość a preferencje przestrzenne}
\label{ch05:top}
\ChapterStatus{draft}{theory-b}

Rozdział rozwija wątek osobowości zapowiedziany w modelu wieloźródłowym preferencji (rozdz.~\ref{ch03:top}). Model Wielkiej Piątki i instrument IPIP-NEO-120 pełnią w nim rolę \emph{hipotezy projektowej} dla syntezy briefu, a nie wyroku estetycznego o „właściwym” guście użytkownika.
Rozdział uzasadnia źródło \texttt{personality} w ślepej macierzy sześciu wariantów artefaktu IDA i wiąże je z pytaniem eksploracyjnym~\RQ{4} \parencite{idaPlatform2026}.
Mapowania domen i facetów na parametry wnętrza dokumentuje załącznik~\ref{app:c:top}, a potok i interfejs macierzy rozdz.~\ref{ch07:top}.

\section{Model Wielkiej Piątki}
\label{ch05:sec:bigfive}

John, Naumann i Soto opisują Wielką Piątkę (\emph{Big Five}) jako taksonomię różnic indywidualnych w pięciu szerokich domenach: otwartość na doświadczenie (\emph{Openness}), sumienność (\emph{Conscientiousness}), ekstrawersja (\emph{Extraversion}), ugodowość (\emph{Agreeableness}) i neurotyczność, czyli emocjonalna labilność (\emph{Neuroticism}). Skrót OCEAN pochodzi od pierwszych liter ich angielskich nazw \parencite{john2008bigfive}.
Model nie wyjaśnia „osobowości jako takiej” w sensie klinicznym. Daje wspólny język cech o względnie stałej strukturze, który przydaje się także poza psychologią różnic indywidualnych, na przykład w projektowaniu interakcji i rekomendacji.

Dla badań przez projektowanie (RtD) ważne jest rozróżnienie: w IDA Big Five nie służy diagnozie uczestnika, tylko operacjonalizuje \emph{jeden z wielu} kanałów preferencji przestrzennych.
Osobowość jest materiałem briefu, podobnie jak preferencje jawne i ukryte (rozdz.~\ref{ch03:top}), i nie rozstrzyga o estetyce.
W ścieżce pełnej (Full) domeny i facety zasilają syntezę promptu. Ścieżka szybka (Fast) nie zawiera kwestionariusza IPIP-NEO-120, więc nie dostarcza danych pod~\RQ{4} (por.\ rozdz.~\ref{ch06:top} i~\ref{ch07:top}).

\section{IPIP i IPIP-NEO-120}
\label{ch05:sec:ipip}

Publiczna pula pozycji IPIP (\emph{International Personality Item Pool}) pozwala budować skale Big Five bez barier licencyjnych, które towarzyszą inwentarzom komercyjnym.
Obecna platforma IDA używa \textbf{IPIP-NEO-120}: 120 pozycji, pięć domen i trzydzieści facetów, po sześć na domenę. Wcześniejsze wersje projektu zakładały formy skrócone (m.in.\ IPIP-60), ale nie opisują stanu obecnego \parencite{idaPlatform2026}.
\NeedsCite{Johnson IPIP-NEO-120 / dokumentacja IPIP; Costa \& McCrae NEO jako tradycja facetów --- weryfikacja DOI przed cytowaniem}

Na poziomie facetów personalizacja ma sens. Sama domena, na przykład wysoka otwartość, jest zbyt szeroka, żeby wskazać styl wnętrza. Facet estetyki ($O2$), fantazji ($O1$) czy porządku ($C2$) pozwala sformułować węższą hipotezę projektową (tab.~\ref{tab:ch05:hipotezy}; zał.~\ref{app:c:top}).
Zapis danych badawczych obejmuje m.in.\ \texttt{big5\_*}, \texttt{big5\_facets}, \texttt{big5\_responses} i znacznik kompletności (\texttt{has\_complete\_bigfive}).
Rozprawa nie powtarza pełnej listy pozycji IPIP. Instrument i scoring opisuje załącznik~\ref{app:b:top}, a tekst główny parafrazuje tylko format pytań, bez treści chronionych pozycji.

\section{Literatura: osobowość a preferencje środowiskowe}
\label{ch05:sec:literatura}

Literatura łącząca cechy osobowości z preferencjami estetycznymi i środowiskami mieszkalnymi wskazuje na asocjacje \emph{słabe do umiarkowanych}, a nie na prawa deterministyczne.
Starsze obserwacje sugerują, że przestrzenie osób o wyższej otwartości bywają bogatsze w zróżnicowane obiekty i elementy artystyczne, a wyższa sumienność idzie w parze z porządkiem i czytelnym układem. Efekt zależy jednak od kontekstu, kultury i metody pomiaru.
\NeedsCite{Gosling i~in.\ 2002 --- A room with a cue (JPSP); weryfikacja DOI}
\NeedsCite{Graham, Gosling \& Travis 2017 --- psychology of home environments; weryfikacja DOI}
\NeedsCite{Nasar 1989 --- symbolic meanings of house styles / Environment and Behavior; weryfikacja DOI}
\NeedsCite{Whitfield \& Wiltshire 1990 --- color psychology (przegląd krytyczny); weryfikacja DOI}

Wniosek projektowy dla IDA jest ostrożny. Literatura uzasadnia \emph{eksploracyjne} mapowanie cech na parametry briefu (złożoność wizualna, materiały, atmosfera), ale nie pozwala twierdzić, że Big Five „wyjaśnia gust”.
Przy małej próbie pilotażowej korelacje domen i facetów z \texttt{implicit\_*}, \texttt{explicit\_*} oraz z wyborem źródła \texttt{personality} mają w \RQ{4} status eksploracyjny. Wartość $r\!>\!0{,}4$ nie może tu służyć jako wynik potwierdzający \parencite{idaPlatform2026}.
Plan analiz i ograniczenia mocy testów opisuje rozdz.~\ref{ch06:top}.

\section{Mapowanie facetów na parametry wnętrz}
\label{ch05:sec:mapowanie}

W potoku IDA wyniki IPIP-NEO-120 są normalizowane i trafiają do reguł stylu (\texttt{BIGFIVE\_STYLE\_MAPPINGS} w module \texttt{prompt-synthesis}).
Warunki na domenach i facetach wskazują dominantę stylu, listę materiałów i złożoność wizualną briefu JSON~v3.
Ta warstwa należy do \emph{explainable personalization}: można sprawdzić, która reguła wpłynęła na źródło \texttt{personality}, a także na \texttt{mixed} i \texttt{mixed\_functional}, gdzie osobowość wchodzi z wagą $0{,}30$.
Szczegóły operacyjne podaje załącznik~\ref{app:c:top}, a scoring i interfejs macierzy rozdz.~\ref{ch07:sec:macierz}.

Tabela~\ref{tab:ch05:hipotezy} zbiera wybrane hipotezy projektowe wraz z ostrzeżeniami interpretacyjnymi. Rozprawa nie dowodzi ich empirycznie.
Są to reguły sterujące syntezą obrazu, które ślepy wybór użytkownika ma sprawdzić.

\begin{table}[htbp]
  \centering
  \caption{Domena / facet IPIP-NEO-120 $\rightarrow$ hipoteza projektowa $\rightarrow$ ostrożność interpretacji
    (skrót mapowań IDA; pełnia w zał.~\ref{app:c:top}).}
  \label{tab:ch05:hipotezy}
  \begin{tabular}{@{}p{0.22\textwidth}p{0.38\textwidth}p{0.30\textwidth}@{}}
    \toprule
    Domena / facet & Hipoteza projektowa (brief) & Ostrożność interpretacji \\
    \midrule
    Otwartość ($O$); estetyka $O2$
      & Wyższa złożoność, eklektyzm, materiały zróżnicowane / artystyczne
      & $O$ obejmuje także intelektualną ciekawość --- nie zawsze „boho” \\
    \addlinespace
    Otwartość ($O$); fantazja $O1$ + niska $C$
      & Styl bardziej swobodny, mniej „czystego” minimalizmu
      & Niska $C$ $\neq$ niechlujstwo; może oznaczać elastyczność funkcji \\
    \addlinespace
    Sumienność ($C$); porządek $C2$
      & Minimalizm, czytelne linie, niska złożoność wizualna
      & Preferencja porządku $\neq$ odrzucenie dekoru; kontekst kulturowy \\
    \addlinespace
    Ekstrawersja ($E$); seeking $E5$
      & Wyższa stymulacja, odważniejsza paleta / wzory
      & Ekstrawersja społeczna $\neq$ maksymalizm estetyczny \\
    \addlinespace
    Ugodowość ($A$) + niska $N$
      & Atmosfera ciepła, harmonijna; miękkie tekstylia
      & Ryzyko stereotypu „przytulnie = ugodowo”; weryfikacja wyborem \\
    \addlinespace
    Neurotyczność ($N$); lęk $N1$
      & Spokojniejsza atmosfera, ograniczona stymulacja
      & Nie patologizować; unikać „bezpieczeństwa” jako piętna cechy \\
    \bottomrule
  \end{tabular}
\end{table}

\NeedsCite{podstawa empiryczna poszczególnych reguł --- m.in.\ tradycja środowisk osobowościowych (Gosling i~in.); nie cytować jako praw przyczynowych bez DOI}

Każda hipoteza spotyka się z pozostałymi kanałami: swipe (implicit), deklaracjami (explicit), funkcją pokoju i inspiracjami.
Źródło \texttt{personality} izoluje w macierzy kanał osobowościowy, a warianty mieszane łączą go z innymi sygnałami. Użytkownik wybiera \emph{obraz}, a nie etykietę cechy (rozdz.~\ref{ch07:sec:slepy-wybor}).

\section{Krytyka i ograniczenia}
\label{ch05:sec:krytyka}

\subsection{Korelacje, nie przyczynowość}
Asocjacje osobowość--estetyka nie uprawniają do wnioskowania przyczynowego ani do „odczytywania” stylu z cechy.
W~\RQ{4} przy małym~$N$ wolno pokazać heatmapy i szerokie przedziały ufności. Nie wolno przedstawiać pilotażu jako walidacji psychometrycznej mapowań.

\subsection{Ryzyko determinizmu i stereotypizacji}
Formuła „jesteś typu~$X$, więc industrialny loft albo skandynawski minimalizm” szkodzi projektowi: odbiera użytkownikowi sprawczość i utrwala stereotypy.
IDA przeciwdziała temu na trzy sposoby: (a)~łączy wiele źródeł, (b)~pokazuje warianty bez etykiet źródła, (c)~traktuje reguły jako hipotezy, które wybór behawioralny może obalić.
Przejrzysty potok nie zwalnia z odpowiedzialności. XAI oznacza tu cytowalne reguły, a nie pełną wyjaśnialność modelu generatywnego. Ma umożliwiać krytykę, a nie uzasadniać determinizm.

\subsection{Etyka, zgoda, RODO --- most do metodologii}
Wnioskowanie o preferencjach z cech osobowości wymaga świadomej zgody, minimalizacji danych i jasno określonego celu przetwarzania.
Procedurę zgód, anonimizację (\texttt{user\_hash}) i RODO opisuje rozdz.~\ref{ch06:sec:etyka} wraz z załącznikiem~\ref{app:f:top}. Tutaj wystarczy zasada: dane Big Five służą personalizacji briefu w badaniu RtD, a nie profilowaniu marketingowemu ani ocenie „jakości” osoby.
\NeedsCite{etyka rekomendacji opartych o osobowość / fair ML --- przegląd peer-review; weryfikacja DOI przed cytowaniem}

W logice RtD artefakt jest argumentem \parencite{zimmerman2007rtd}. Źródło \texttt{personality} sprawdza, czy kanał cech wnosi cokolwiek poza kanałami implicit i explicit, a odpowiedź daje wybór w macierzy, nie deklaracja badacza.

\section{Uzasadnienie włączenia osobowości do pipeline IDA}
\label{ch05:sec:uzasadnienie}

Osobowość wchodzi do IDA jako trzeci filar estetyczny obok preferencji ukrytych i jawnych (rozdz.~\ref{ch03:sec:synteza}):
\begin{enumerate}
  \item \textbf{Źródło \texttt{personality}}: brief powstaje wyłącznie z mapowań IPIP-NEO-120, a pozostałe kanały estetyczne są wyzerowane.
  \item \textbf{Wkład do \texttt{mixed} i \texttt{mixed\_functional}}: kanoniczna waga wynosi $0{,}30$ (obok implicit $0{,}40$ i explicit $0{,}30$). Funkcja pokoju dochodzi dopiero w \texttt{mixed\_functional}.
  \item \textbf{Pytanie~\RQ{4}}: czy domeny i facety wiążą się z parametrami profilu (\texttt{implicit\_*}, \texttt{explicit\_*}) oraz z preferencją wariantu osobowościowego w ślepym wyborze (hipotezy robocze H4a--c, status eksploracyjny przy małym~$N$).
\end{enumerate}

Big Five trafia do IDA nie dlatego, że determinizm estetyczny jest prawdopodobny, lecz ze względu na \emph{kontrast eksperymentalny}. Bez osobnego wariantu nie da się sprawdzić, czy użytkownicy w ogóle wybierają koncepcje wyprowadzone z cech, gdy nie znają etykiety źródła (\RQ{3}/\RQ{4}).
Osobowość pozostaje wątkiem pomocniczym wobec głównej osi empirycznej, czyli rozbieżności implicit i explicit oraz wyborów w macierzy. Zgadza się to z pozycją tej rozprawy w dyscyplinie sztuk, w której Big Five nie może przejąć narracji projektowej.

IPIP-NEO-120 jest więc w IDA instrumentem hipotez projektowych, cytowalnym wejściem przejrzystego potoku i warunkiem eksploracji~\RQ{4}. Nie jest wyrokiem o guście ani zamiennikiem dialogu z użytkownikiem.
